\chapter*{まえがき}\label{maegaki}\addcontentsline{toc}{chapter}{まえがき}\markboth{まえがき}{まえがき}

皆さん、はじめまして！この本は「構文解析」というテーマについて扱った、一風変わった本です。これまで「構文解析」については、コンパイラや言語処理系を扱った書籍の一部で触れられる程度でした。「Parsing
Techniques」という英語の書籍があるものの、英語圏ですら広く知られているとは言い難いですし、内容も網羅的ではあるものの平易とは言えません。

私がはじめて「構文解析」の入口に立ったのは高校三年生のときでした。当時、私はプログラミング言語に興味を持ち始め、いつかは自分のプログラミング言語を作ってみたいと思っていました。しかし、高校の図書館にあった「プログラミング言語を作る」本を借りて読んでみたものの、とても難しく当時の私には理解できませんでしたし、コンパイラ開発者のバイブルとも呼ばれる「ドラゴンブック」\footnote{Alfred
  V. Aho, Monica S. Lam, Ravi Sethi, Jeffrey D. Ullman 著「コンパイラ
  第2版: 原理・技法・ツール」（サイエンス社,
  2009年）のこと。コンパイラ開発の標準的な教科書であり、表紙のデザインから「ドラゴンブック」の愛称で広く知られています。}も同様でした。

何がきっかけかは覚えていないのですが、いきなり大きな言語を作ろうとしても難しいということでまずは括弧を含む算術式を計算できる「電卓」アプリを作ることにしました。非常に単純なもので

\begin{codescreen}{}
(1 + 2) * (3 + 4)
\end{codescreen}

\noindent のようにテキストフィールドに入力して「計算」ボタンを押すと\texttt{21}と表示される、ただそれだけの代物です。今振り返れば、この電卓アプリはごく初歩的な「構文解析」と「インタプリタ」を実装したものとも言えます。

ただし、当時の私は言語処理系を作る際に必須の知識がほとんどなく、抽象構文木も典型的な構文解析のアルゴリズムも知らず、足りない知恵を振り絞って括弧の対応や演算子の優先順位を自力で計算したのでした。それは本当に拙いものでしたが、それでもなんとか動くものを作ることができたのです。

それから二十年余り。筆者は現在、研究職ではないものの構文解析の手法の一つである、Parsing
Expression
Grammar（PEG）を専門として、時折関連する論文の査読を引き受けたりしています。

現在の私は研究に未練を残しつつちょっとアカデミアに関わりがある微妙な立場ですが、そんな中でも折に触れて思うことがありました。構文解析はなぜ、一般の技術書であまり取り上げられないのだろうか、ということです。

言語処理系を扱う書籍の一部として構文解析が取り上げられることは珍しくありません。しかし、その扱いはあくまで「おまけ」であって、言語を作る上での通過点としてしか位置づけられていません。これについては、言語の本質は構文解析の「後」にあるのであり、構文は本質ではないというのが大きいでしょう。筆者もこの点に異論はありません。ただ、構文（正確には具象構文）はほんとうに「おまけ」なのでしょうか。ときどき疑問に思います。

プログラミング言語に携わる研究者の間では、プログラミング言語の「意味」は抽象構文木に対して定められるものであり、構文は「ガワ」に過ぎないという考え方の人が多いように思いますし、その考えも間違ってはいないと思います。

しかし、現実に我々が読み書きするのはプログラムの抽象構文木ではなく、プログラムの文字列です。構文はUIであると言えます。多くのアプリケーションにおいてUIが果たす役割は非常に大きく、UI自体が専門分野として存在しているくらいです。そのUIたる構文とUIを適切な内部構造に変換する構文解析は軽視されていいのだろうか。筆者が構文解析の本を書こうと思い立ったのはそんな「こだわり」からでした。

本書の読者の皆さんには、構文解析の世界を少しでも楽しんでいただければと思います。構文解析は非常に奥深いテーマであり、その奥深さを一冊の本で完全に網羅することはできません。しかし、本書を通じて構文解析の基礎を学び、その面白さを感じていただければと思います。

\section*{本書のアプローチ}\markboth{まえがき}{本書のアプローチ}

本書は、理論から入って応用に降りてくる、という順序をとりません。逆です。まず小さな構文解析器を手で書き、動かし、なぜ動くのかを考え、そのあとで理論に戻ります。

この順序を選んだ理由は単純で、構文解析は手を動かさないとわからないからです。オートマトンの状態遷移図を眺めているだけでは、なぜ括弧の対応が数えられるのかは腑に落ちません。しかし、自分で書いた再帰下降の構文解析器が括弧を正しく処理するのを見れば、その瞬間に理解が追いつきます。

本書のコードが、特別な言語機能に頼らない素直な書き方で統一されているのも同じ理由によります。読者が普段使っている言語へ移し替えながら読んでいただいても構いません。

\section*{本書が想定する読者}\markboth{まえがき}{本書が想定する読者}

本書は、何らかのプログラミング言語で日常的にコードを書いている方を想定しています。サンプルコードは主にJavaで書かれていますが、Javaに習熟している必要はありません。クラスとメソッドが読め、条件分岐と繰り返しが追えれば十分です。

前提知識として、コンパイラや形式言語理論をあらかじめ学んでおく必要はありません。必要な概念は、それを使う場面でその都度説明します。数学的な厳密さよりも「なぜそう考えるのか」を優先しました。

一方で、本書は入門書であると同時に、入門で終わらない本でもあります。電卓やJSONの構文解析器を書き上げたあとには、LL法やLR法といった古典的なアルゴリズムの内側に踏み込み、さらに実在のプログラミング言語が抱える厄介な問題まで扱います。読み進めるにつれて話は確実に難しくなりますが、その難しさこそが構文解析の面白さでもあります。

\section*{本書の構成}\markboth{まえがき}{本書の構成}

本書は次のような構成になっています。

序章では、構文解析がたどってきた歴史を概観し、この技術がそもそも何をするものなのか、そしてどれほど身近なところで働いているのかを見ていきます。本書全体の見取り図にあたる部分です。

第1章では簡単な算術式を例題にして、構文解析とはどういう処理かを体感してもらいます。定義を天下り式に提示するような本もありますが、構文解析については「まずは書いてみる」のが手っ取り早いというのが筆者の持論です。

第2章ではJSONの構文解析器を書いてもらいます。JSONは世界中で使われている非常に実用的なデータ形式（言語）でありながら、非常にシンプルです。そのシンプルなJSONを通して、実用的な言語の構文解析の基礎を学んでいただければと思います。

第3章では文脈自由文法（や言語）について解説します。文脈自由文法やそれに関する理論は構文解析の基礎になっています。括弧の対応を表現する言語を\texttt{Dyck}言語と言いますが、\texttt{Dyck}言語は文脈自由言語を特徴づけるものです。この章を理解することで文脈自由文法の直観的な理解が得られます。現代の構文解析は文脈自由文法を基盤としていますから、文脈自由言語の概念を理解することは非常に重要です。

第4章では現在の構文解析で広く採用されているアルゴリズムのうち、主だったものを解説します。特にLL(1)やLR(1)といったアルゴリズムについてできるだけ平易にかつ詳しく説明します。また、PEGやPackrat
Parsingという比較的新しい構文解析アルゴリズムについても詳しく説明します。

第5章では構文解析器生成系について解説します。Yaccのような古典的なものに留まらず、ALL(*)アルゴリズムを基盤にしたANTLRやPEGを基盤にしたパーサーコンビネータなど、最新の構文解析器生成系について説明します。さらに、第5章では簡単なパーサーコンビネータを自作します。パーサーコンビネータは元々は関数型プログラミング言語から出てきたテクニックですが昨今ではいろいろな言語でパーサーコンビネータライブラリがあります。パーサーコンビネータを自作する体験を通じて「文の定義から構文解析器を生成する」とはどういうことか理解してもらえるのではないかと思います。

第6章では従来の言語処理系についての本が取り扱わなかった「現実の構文解析」の話をします。従来の書籍に書かれている構文解析の世界はとても「きれい」なものです。しかし、RubyでもPythonでもあるいはJavaScriptでもよいですが、現実の構文解析は必ずしもそこまできれいにはいかないものです。この章を通して現実の言語における構文解析はとても泥臭いものであること、その泥臭さを通して「書きやすく読みやすい」文法が実現されていることを実感してもらえるのではないかと思います。

「おわりに」は締めくくりとしてこれまでの章を振り返りつつ、今後、皆さんが構文解析を学ぶにあたって参考になりそうな文献や資料について紹介します。この本は構文解析のみを取り扱った珍しい本ですが、それでも構文解析の世界は広く、本書で取り扱わなかったテーマも多々あります。この章を読んで、構文解析の世界により深く興味を持っていただければ幸いです。

\section*{本書が扱わないこと}\markboth{まえがき}{本書が扱わないこと}

範囲をはっきりさせるために、本書が扱わないことも述べておきます。

まず、自然言語の解析は対象外です。本書でいう構文解析は、プログラミング言語やデータ形式のように、人間が設計した厳密な文法を持つテキストを相手にします。日本語や英語のように曖昧さを本質的に抱えた言語の解析は、まったく別の技術体系に属します。

次に、構文解析より後ろの工程、すなわち意味解析、型検査、最適化、コード生成には踏み込みません。いずれもそれぞれで一冊が書ける主題です。入力から抽象構文木を組み立てるところまでが、本書の守備範囲になります。

また、構文解析器の速度を極限まで高めるための技法も扱いません。本書のコードは、速さよりも読んでわかることを優先して書かれています。実務で性能が問題になる場面は確かにありますが、それは仕組みを理解したあとの話でしょう。

\section*{本書で扱うプログラミング言語とツール}\markboth{まえがき}{本書で扱うプログラミング言語とツール}

本書では、構文解析を学ぶために、以下のプログラミング言語とツールを使用します。

\textbf{\sffamily Java}:
本書の主要な実装言語としてJavaを選びました。Javaは静的型付けであり、抽象構文木のような複雑なデータ構造を表現するのにも不足はありません。本書ではJava
21を基準に解説します。なお、最新のLTSであるJava
25（2025年9月リリース）でも本書のコードはそのまま動作します。

\textgt{パーサージェネレータ}: 第5章では以下のツールを扱います。

\begin{itemize}
\item
  ANTLR4：現代的で強力なパーサージェネレータ
\item
  JavaCC：Java専用の伝統的なパーサージェネレータ
\item
  Yacc/Lex：C言語用の古典的ツール（歴史的理解のため）
\end{itemize}

\textgt{開発環境}: 特定のIDEは必須ではありませんが、IntelliJ IDEA
Community Edition（以下、IDEA）を推奨します。IDEAはJavaの開発に適しており、構文解析器の実装やテストを効率的に行うことができます。Visual Studio Codeのようなエディタでも問題ありません。

すべてのサンプルコードは、下記のGitHubリポジトリから入手可能です。

\begin{quote}
\url{https://github.com/kahei/parser_book}
\end{quote}

サンプルコードは、このリポジトリの\texttt{code}ディレクトリ以下に、章ごとの独立したプロジェクトとして置かれています。本文中のコードは紙面の都合で一部を省略している箇所がありますので、全体を動かしたい場合はこちらを参照してください。

\section*{学習の道筋}\markboth{まえがき}{学習の道筋}

本書はさまざまな背景を持つ読者を想定しています。ここでは、目的に応じた読み進め方を提案します。

\noindent\textgt{初心者の方へ}

\begin{itemize}
\item
  まず第1章で構文解析の基本的な考え方を理解してください
\item
  次に第2章でJSONパーサーを実装し、実践的な経験を積みます
\item
  その後、第5章でツールを使った構文解析器の生成を学びます
\item
  理論的な内容（第3章、第4章）は後回しにしても構いません
\end{itemize}

\noindent\textgt{理論もしっかり学びたい方へ}

\begin{itemize}
\item
  第1章→第2章→第3章→第4章の順番で読み進めてください
\item
  第3章の文脈自由文法は構文解析の理論的基礎です
\item
  第4章のアルゴリズムは少し難しいですが、じっくり取り組んでください
\item
  第5章で理論の実装を、第6章で実践を学びます
\end{itemize}

\noindent\textgt{すぐに使える知識が欲しい方へ}

\begin{itemize}
\item
  第1章→第2章→第5章→第6章の順で読むことをお勧めします
\item
  第6章では現実に遭遇する泥臭い構文解析の問題を扱います
\end{itemize}

\noindent\textgt{各章の難易度の目安}

\begin{itemize}
\item
  第1章：★（入門）
\item
  第2章：★★（実践基礎）
\item
  第3章：★★★（理論基礎）
\item
  第4章：★★★★（理論応用）
\item
  第5章：★★（ツール活用）
\item
  第6章：★★（実践応用）
\item
  おわりに：★（まとめ）
\end{itemize}

最初から最後まで読み通してもらえれば、最終的には構文解析の基礎から応用まで幅広く理解できるようになります。

\section*{表記について}\markboth{まえがき}{表記について}

本書では、術語は原則として日本語で表記します。「構文解析器」（パーサー）、「構文解析器生成系」（パーサージェネレータ）、「抽象構文木」（AST）などがそれにあたります。カタカナ語や略語のほうが通りがよい場面もありますが、はじめて学ぶ方にとっては、漢字のほうが意味を推測しやすいと考えたためです。

ただし、「パーサーコンビネータ」のようにカタカナでの呼び名が定着しているものは、そのまま用います。

\section*{謝辞}\markboth{まえがき}{謝辞}

本書は、多くの方々の力を借りて形になりました。

草稿の段階で原稿に目を通し、貴重なご指摘をくださった以下の方々に、心よりお礼申し上げます（敬称略、五十音順）。専門的な観点からのご指摘によって、説明の足りない箇所が補われ、記述の正確さが幾度となく改善されました。

\begin{itemize}
\item
  市川和央
\item
  きしだなおき（kis）
\item
  倉光君郎（日本女子大学）
\item
  佐藤貴比呂（株式会社kubell）
\item
  なぎせゆうき
\end{itemize}

もちろん、本書に残る誤りはすべて筆者の責任です。

ZEN大学出版会編集長の鈴木嘉平さんにも大変お世話になりました。私が最初に企画を持ちかけたのが2020年頃ですから、実に6年越しの企画となりました。毎月の定例ミーティングで進捗を報告できない月もありましたが、根気強く待っていただきました。

最後に、妻の久美子には本当に助けられました。執筆の途中、何度かめげそうになったことがありましたが、そのたびに励ましてくれました。

\bigskip

\noindent
2026年5月、自室にて記す。
